\documentclass[UTF8,a4paper,12pt]{ctexart}

\usepackage[a4paper,margin=2.4cm]{geometry}
\usepackage{amsmath,amssymb}
\usepackage{array,booktabs,longtable,tabularx}
\usepackage{enumitem}
\usepackage{xcolor}
\usepackage{hyperref}
\usepackage{fancyhdr}

\hypersetup{
  colorlinks=true,
  linkcolor=blue,
  urlcolor=blue,
  pdftitle={四因子计算交付说明},
  pdfauthor={}
}

\setlength{\parindent}{2em}
\setlength{\parskip}{0.35em}
\setlist[itemize]{leftmargin=2.5em}
\setlist[enumerate]{leftmargin=2.8em}
\renewcommand{\arraystretch}{1.25}
\newcolumntype{Y}{>{\raggedright\arraybackslash}X}
\newcommand{\pathtext}[1]{\texttt{\detokenize{#1}}}
\newcommand{\code}[1]{\texttt{\detokenize{#1}}}

\pagestyle{fancy}
\fancyhf{}
\lhead{四因子计算交付说明}
\rhead{\thepage}

\begin{document}

\begin{titlepage}
\centering
\vspace*{3cm}
{\Huge\bfseries 四因子计算交付说明\par}
\vspace{1.2cm}
{\Large EP、PB、ROE 与 Size\par}
\vspace{2cm}
\begin{tabular}{rl}
数据区间：& 2020-01-02 至 2025-12-31\\
时间口径：& T 日因子用于 T+1 日交易\\
输出格式：& date + stock\_code 索引，signal 列\\
\end{tabular}
\vfill
{\large 版本：汇总交付版\par}
{\large \today\par}
\end{titlepage}

\tableofcontents
\newpage

\section{项目范围}

本文档汇总 EP、PB、ROE、Size 四个股票横截面因子的计算流程、数据清洗规则、点时间对齐方式、结果格式和验证结果。

统一时间口径为：

\begin{quote}
T 日收盘后使用截至 T 日结束时已经可获得的数据计算因子；T 日因子用于 T+1 日交易。
\end{quote}

本项目完成的是因子计算和数据验证，不包含 T+1 实际成交、涨跌停、成交量限制、滑点、交易成本和完整回测。

\section{交付目录和正式结果}

\begin{table}[htbp]
\centering
\caption{四个因子的交付目录}
\begin{tabularx}{\textwidth}{lYY}
\toprule
因子 & 交付目录 & 正式结果\\
\midrule
EP & \pathtext{D:\因子计算\市盈率\EP因子_交付版} & \pathtext{因子结果\ep.parquet}\\
PB & \pathtext{D:\因子计算\市净率\PB因子_交付版} & \pathtext{因子结果\pb.parquet}\\
ROE & \pathtext{D:\因子计算\ROE因子_交付版} & \pathtext{因子结果\roe.parquet}\\
Size & \pathtext{D:\因子计算\规模因子_交付版} & \pathtext{因子结果\size.parquet}\\
\bottomrule
\end{tabularx}
\end{table}

四个正式结果均采用同一格式：

\begin{verbatim}
index = ['date', 'stock_code']
columns = ['signal']
\end{verbatim}

其中 \code{signal} 是经过每日横截面 MAD 去极值和 Z 标准化后的因子值。

\section{输入数据}

\subsection{行情数据}

原始行情表为：

\pathtext{D:\实习生学习项目\基础数据\chn_equ_mkt_quotation.parquet}

主要字段包括：

\begin{itemize}
  \item \code{date}；
  \item \code{stock_code} 或 \code{code6}；
  \item \code{me_total}；
  \item \code{suspended}；
  \item \code{status}。
\end{itemize}

四个因子统一复用已经清洗的行情表：

\pathtext{D:\因子计算\市盈率\EP因子_交付版\中间结果\_mkt_clean.parquet}

行情清洗内容如下：

\begin{enumerate}
  \item 股票代码统一为六位代码；
  \item 检查 \code{date + stock\_code} 是否重复；
  \item 保留总市值 \code{me\_total}；
  \item 使用 \code{suspended} 和 \code{status} 形成停牌掩码；
  \item 使用行情表中的首次出现日期识别次新股；
  \item 由于没有完整的历史 ST/PT 标记，当前没有可靠完成历史 ST/PT 过滤；
  \item 由于行情数据从 2020-01-02 开始，首次出现日期只能识别数据集开始后新出现的股票，不能替代官方上市日期。
\end{enumerate}

\subsection{利润表}

原始利润表为：

\pathtext{D:\实习生学习项目\基础数据\vw_fdmt_is_new.parquet}

使用字段为：

\begin{verbatim}
TICKER_SYMBOL, ACT_PUBTIME, UPDATE_TIME, END_DATE,
FISCAL_PERIOD, MERGED_FLAG, N_INCOME_ATTR_P
\end{verbatim}

\code{N\_INCOME\_ATTR\_P} 是归属于母公司股东的净利润。

\subsection{资产负债表}

原始资产负债表为：

\pathtext{D:\实习生学习项目\基础数据\vw_fdmt_bs_new.parquet}

使用字段为：

\begin{verbatim}
TICKER_SYMBOL, ACT_PUBTIME, UPDATE_TIME, END_DATE,
FISCAL_PERIOD, MERGED_FLAG, T_EQUITY_ATTR_P
\end{verbatim}

\code{T\_EQUITY\_ATTR\_P} 是归属于母公司股东权益。

\section{财报清洗和版本处理}

\subsection{报表口径}

利润表和资产负债表均保留合并报表：

\[
\texttt{MERGED\_FLAG}=1.
\]

只保留累计报告期，要求报告期月份等于 \code{FISCAL\_PERIOD}，即 3 月、6 月、9 月和 12 月。

利润表使用累计归母净利润计算 TTM；资产负债表使用报告期末归母股东权益，不计算 TTM。

\subsection{重复记录与修订版本}

完全相同的重复记录只保留一条。

同一股票、同一报告期存在不同财报版本时，不能简单使用第一条记录，也不能把最新版本回填到历史日期。所有版本均保留，并根据版本可用时间参与点时间对齐。

\subsection{财报可用时间}

定义财报版本可用时间：

\begin{equation}
\mathrm{VERSION\_TIME}
=
\max(\mathrm{ACT\_PUBTIME},\mathrm{UPDATE\_TIME}).
\end{equation}

T 日只能使用：

\begin{equation}
\mathrm{VERSION\_TIME}
\leq
\text{T 日 }23{:}59{:}59.999999999.
\end{equation}

因此，T 日更新或公布的财报可以用于计算 T 日因子，并用于 T+1 日交易。后续公告或修订不能回填到修订发生之前的日期。

\subsection{权益表异常记录}

权益表清洗统计如下：

\begin{table}[htbp]
\centering
\caption{资产负债表清洗统计}
\begin{tabular}{lr}
\toprule
项目 & 数量\\
\midrule
原始记录 & 453,920\\
\code{MERGED\_FLAG=1} & 272,641\\
\code{T\_EQUITY\_ATTR\_P} 缺失 & 2\\
无法解析为六位股票代码 & 10,793\\
完全重复额外删除 & 2\\
最终权益版本 & 266,908\\
负归母权益 & 1,897\\
权益等于零 & 0\\
有效版本时间冲突 & 0\\
\bottomrule
\end{tabular}
\end{table}

处理原则如下：

\begin{enumerate}
  \item \code{T\_EQUITY\_ATTR\_P} 缺失记录不进入权益版本表，不填 0，也不用前后期数据填补；
  \item A19110、A20471 等无法可靠解析为六位证券代码的记录不参与因子计算，不直接截取数字或补零匹配；
  \item 原始异常记录仅被排除出计算链路，不代表原始文件被删除；
  \item 同一报告期不同修订版本保留；
  \item 同一股票内部完全相同的重复记录只保留一条；不同股票之间出现相同权益值不属于重复。
\end{enumerate}

\section{EP 因子}

\subsection{计算流程}

EP 的计算步骤为：

\begin{enumerate}
  \item 清洗利润表并保留合并报表、累计口径；
  \item 按 \code{VERSION\_TIME} 保留财报的历史版本；
  \item 计算点时 TTM 归母净利润；
  \item 将 TTM 利润与 T 日总市值进行点时间对齐；
  \item 计算原始 EP；
  \item 进行每日横截面 MAD 去极值和 Z 标准化。
\end{enumerate}

年度报告的 TTM 为当年累计归母净利润。非年度报告采用：

\begin{equation}
\mathrm{TTM}
=
\text{本期累计利润}
+
\text{上年年报利润}
-
\text{上年同期累计利润}.
\end{equation}

EP 公式为：

\begin{equation}
\mathrm{EP}_T
=
\frac{\text{T 日 TTM 归母净利润}}
{\text{T 日总市值}}.
\end{equation}

负利润保留，因此负 EP 保留；利润为 0 时 EP 为 0。总市值小于等于 0、利润缺失、停牌或次新股掩码命中时，EP 为空。

\subsection{实际数据案例}

股票 000002 的 2022-03-31 归母净利润存在两个版本。2023 年修订前使用第一个版本，修订后才使用第二个版本，不把第二个版本回填到 2022 年。

利润表中还存在 85 个相同股票、报告期和 \code{ACT\_PUBTIME}、但 \code{N\_INCOME\_ATTR\_P} 不同的冲突键，共 170 行。处理时使用 \code{UPDATE\_TIME} 形成 \code{VERSION\_TIME}，不随机使用第一行，也不使用 \code{groupby.first()} 拼接字段。

例如，股票 000029 在 2020-01-02 的 \code{suspended=1}，因此该日 EP 为缺失值。如果股票 T 日停牌、T+1 日复牌，严格 T 日因子口径下，T 日仍没有因子；复牌后重新计算的因子用于之后的交易日。

\subsection{结果统计}

\begin{table}[htbp]
\centering
\caption{EP 结果统计}
\begin{tabular}{lr}
\toprule
项目 & 数量\\
\midrule
数据日期 & 2020-01-02 至 2025-12-31\\
行情股票数 & 6,016\\
至少有一次有效 EP & 5,561 只\\
原始有效 EP 记录 & 6,303,328\\
负 EP & 1,422,534\\
MAD 去极值记录 & 701,329\\
最终有效 Z EP & 6,303,327\\
未来数据违规 & 0\\
重复索引 & 0\\
\bottomrule
\end{tabular}
\end{table}

\section{PB 因子}

\subsection{计算流程}

PB 的计算步骤为：

\begin{enumerate}
  \item 复用清洗后的行情表；
  \item 清洗资产负债表并保留合并报表、累计报告期；
  \item 按 \code{VERSION\_TIME} 形成点时归母股东权益快照；
  \item 将最新可用权益与 T 日总市值对齐；
  \item 计算原始 PB；
  \item 进行每日横截面 MAD 去极值和 Z 标准化。
\end{enumerate}

PB 公式为：

\begin{equation}
\mathrm{PB}_T
=
\frac{\text{T 日总市值}}
{\text{T 日最新可获得的归母股东权益}}.
\end{equation}

归母权益小于 0 时保留负 PB。负 PB 表示负净资产风险，不能简单解释为低估值，但保留其信息和排序结果。归母权益等于 0 时 PB 无定义，设为缺失；总市值小于等于 0 或权益缺失时，PB 也设为缺失。

\subsection{实际数据案例}

股票 002499 的 2018-12-31 报告和股票 603895 的 2023-12-31 报告各自存在完全重复入库记录。这些记录属于同一股票内部重复，不是两只股票互相重复。相同股票、报告期、公告时间、更新时间和权益值完全相同的重复行只保留一条。

股票 000023 的部分报告期归母权益为负，对应 PB 保留为负值，并在掩码中标记负净资产。

\subsection{结果统计}

\begin{table}[htbp]
\centering
\caption{PB 结果统计}
\begin{tabular}{lr}
\toprule
项目 & 数量\\
\midrule
行情总行数 & 7,155,034\\
行情股票数 & 6,016\\
资产负债表最终版本 & 266,908\\
多版本报告期 & 37,760\\
原始有效 PB & 6,325,068\\
负 PB & 47,043\\
MAD 去极值 & 580,242\\
最终有效 Z PB & 6,325,067\\
未来数据违规 & 0\\
重复索引 & 0\\
验证结果 & \code{verify\_pb.py PASS}\\
\bottomrule
\end{tabular}
\end{table}

\section{ROE 因子}

\subsection{数据复用}

ROE 不重复清洗利润表，而是直接复用 EP 已经完成的利润版本和 TTM 结果：

\begin{itemize}
  \item \pathtext{D:\因子计算\市盈率\EP因子_交付版\中间结果\reports.parquet}；
  \item \pathtext{D:\因子计算\市盈率\EP因子_交付版\中间结果\reports_ttm.parquet}。
\end{itemize}

不能直接复用 \code{ep.parquet}，因为该文件已经是标准化后的 EP 信号。ROE 使用的是 EP 的中间利润数据。

ROE 同时使用 PB 清洗得到的 \code{T\_EQUITY\_ATTR\_P}。利润和权益均采用合并报表中的归母口径，归属范围一致。

\subsection{计算公式}

最新报告期权益和去年同期权益分别点时对齐：

\begin{equation}
\text{平均归母权益}
=
\frac{\text{本期归母权益}+\text{去年同期归母权益}}{2}.
\end{equation}

ROE 公式为：

\begin{equation}
\mathrm{ROE}_T
=
\frac{\text{TTM 归母净利润}}
{\text{平均归母股东权益}}.
\end{equation}

当前处理规则如下：

\begin{itemize}
  \item 平均权益本期或去年同期缺失时，ROE 为空；
  \item 平均权益小于等于 0 时，ROE 为空，并在掩码中记录负权益；
  \item 负净利润保留，因此可以产生负 ROE；
  \item 停牌和次新股掩码命中时，不进入有效 ROE 样本。
\end{itemize}

\subsection{实际数据案例}

2022-05-05 股票 000002 的独立复算结果如下：

\begin{table}[htbp]
\centering
\caption{ROE 独立复算}
\begin{tabular}{lr}
\toprule
项目 & 数值\\
\midrule
TTM 归母净利润 & 22,660,741,167.06\\
平均归母权益 & 231,289,043,606.325\\
手工 ROE & 0.0979758522657504\\
原始 ROE 输出 & 0.0979758522657504\\
绝对误差 & 0\\
\bottomrule
\end{tabular}
\end{table}

\subsection{结果统计}

\begin{table}[htbp]
\centering
\caption{ROE 结果统计}
\begin{tabular}{lr}
\toprule
项目 & 数量\\
\midrule
行情行数 & 7,155,034\\
EP 利润表版本 & 266,865\\
EP TTM 事件 & 131,134\\
同时有利润和权益数据的股票 & 6,461\\
ROE 事件 & 257,149\\
原始有效 ROE & 5,361,049\\
负 ROE & 1,207,621\\
MAD 去极值 & 496,180\\
横截面标准差为 0 的日期 & 1\\
最终有效 Z ROE & 5,361,048\\
未来数据违规 & 0\\
验证结果 & \code{verify\_roe.py PASS}\\
\bottomrule
\end{tabular}
\end{table}

\section{Size 因子}

\subsection{数据来源和计算公式}

Size 只需要清洗后的行情表，不使用利润表、资产负债表或 TTM。计算字段为 \code{date}、\code{stock\_code}、\code{me\_total}、\code{suspended} 和 \code{is\_cixin}。

原始规模因子采用总市值的自然对数：

\begin{equation}
\mathrm{Size}_{\mathrm{raw},T}
=
\ln(\text{T 日总市值}).
\end{equation}

计算要求如下：

\begin{itemize}
  \item \code{me\_total} 缺失或小于等于 0 时，Size 为空；
  \item T 日停牌或次新股掩码命中时，Size 为空；
  \item T 日 Size 用于 T+1 日交易；
  \item T 日停牌、T+1 日复牌时，T 日仍无新 Size，复牌后的有效因子用于之后的交易日。
\end{itemize}

\code{signal} 越高表示市值越大，越低表示市值越小。如果构造小市值方向组合，回测时可以使用 \(-\code{signal}\)。

\subsection{与 Fama--French SMB 的区别}

本项目的 Size 是逐股票规模暴露：

\begin{verbatim}
index = ['date', 'stock_code']
columns = ['signal']
\end{verbatim}

它不是 Fama--French 的 SMB 组合收益。SMB 需要按市值将股票分为小市值组和大市值组，再计算两组未来收益之差，输出的是组合时间序列。

\section{MAD 去极值和 Z 标准化}

四个因子使用相同的每日横截面处理。

对每个交易日的有效原始因子计算：

\begin{align}
\mathrm{median} &= \text{当日横截面中位数},\\
\mathrm{MAD} &= \operatorname{median}(|x-\mathrm{median}|),\\
\mathrm{scale} &= 1.4826\times\mathrm{MAD},\\
\mathrm{lower} &= \mathrm{median}-3\times\mathrm{scale},\\
\mathrm{upper} &= \mathrm{median}+3\times\mathrm{scale}.
\end{align}

首先进行截断：

\begin{equation}
x_{\mathrm{MAD}}
=
\min\left(\max(x,\mathrm{lower}),\mathrm{upper}\right).
\end{equation}

然后进行横截面 Z 标准化：

\begin{equation}
Z
=
\frac{x_{\mathrm{MAD}}-\overline{x}_{\mathrm{MAD}}}
{\sigma_{\mathrm{MAD}}}.
\end{equation}

处理规则如下：

\begin{enumerate}
  \item 只使用当日有效股票横截面计算中位数、MAD、均值和标准差；
  \item MAD 为 0 时不进行截断；
  \item 横截面标准差为 0 时，Z 值设为缺失，不填 0；
  \item 负 EP、负 PB 和负 ROE 按各自口径参与处理，不事先删除；
  \item 标准化后的 \code{signal} 不代表原始因子的经济单位。
\end{enumerate}

2020-01-02 有效样本数量很少，横截面标准差为 0，因此该日 Z 值为空，这是数学上的正常结果，不是程序错误。

\section{停牌和回测处理边界}

当前因子文件只按 T 日行情状态形成因子掩码，未模拟实际成交。建议回测采用以下规则：

\begin{table}[htbp]
\centering
\caption{不同停牌期限的建议处理}
\begin{tabularx}{\textwidth}{l l Y Y Y}
\toprule
类型 & 连续停牌交易日 & 新建仓 & 已持仓 & 复牌处理\\
\midrule
短期 & 1--5 天 & 停牌期间不买入 & 保留，按最后成交价估值 & 取得有效收盘价后重新计算\\
中期 & 6--20 天 & 停牌期间不买入 & 保留并标记不可交易 & 复牌日观察，下一交易日再考虑新仓\\
长期 & 超过 20 天 & 从候选池剔除 & 保留，不能强制卖出 & 复牌后观察 3 个交易日，再考虑新仓\\
\bottomrule
\end{tabularx}
\end{table}

停牌期间不能把持仓股票删除，也不能将缺失价格填成 0。一般沿用最后成交价估值；复牌后的价格跳变应在复牌时计入组合收益和最大回撤。

在严格 T 日因子、T+1 交易口径下：

\begin{itemize}
  \item T 日停牌导致因子缺失时，T+1 不新买入；
  \item T+1 复牌后重新取得有效行情并计算因子；
  \item 新因子最早用于 T+2 交易；
  \item 不建议将停牌前旧因子无限期向后填充。
\end{itemize}

\section{验证内容和结果}

四个因子均检查以下内容：

\begin{enumerate}
  \item 输出索引是否为 \code{date}、\code{stock\_code}；
  \item 输出列是否只有 \code{signal}；
  \item \code{date + stock\_code} 是否重复；
  \item \code{signal} 是否存在无穷值；
  \item 原始公式是否与中间数据一致；
  \item MAD 截断记录数是否与统计表一致；
  \item 每日 Z 值均值是否接近 0；
  \item 每日 Z 值标准差是否接近 1；
  \item 点时间数据是否满足 \code{VERSION\_TIME} 不晚于 T 日；
  \item 因子是否存在未来数据泄漏。
\end{enumerate}

验证结果为：

\begin{table}[htbp]
\centering
\caption{四因子验证结果}
\begin{tabular}{lll}
\toprule
因子 & 验证脚本 & 结果\\
\midrule
EP & \code{factor\_verify.py} & 通过，未来数据违规 0\\
PB & \code{verify\_pb.py} & PASS，未来数据违规 0\\
ROE & \code{verify\_roe.py} & PASS，未来数据违规 0\\
Size & \code{verify\_size.py} & 检查通过，未来数据违规 0\\
\bottomrule
\end{tabular}
\end{table}

\section{运行顺序}

\subsection{EP}

\begin{verbatim}
python step1_样本筛选.py
python step2_利润表处理.py
python step3_TTM计算.py
python step4_点时间对齐_计算PE.py
python step5_MAD去极值_Z标准化.py
python factor_verify.py
\end{verbatim}

\subsection{PB}

\begin{verbatim}
python pb_step1_读取已清洗行情.py
python pb_step2_资产负债表处理.py
python pb_step3_点时权益对齐.py
python pb_step4_点时间对齐_计算PB.py
python pb_step5_MAD去极值_Z标准化.py
python verify_pb.py
\end{verbatim}

\subsection{ROE}

\begin{verbatim}
python roe_step1_读取已清洗行情.py
python roe_step2_利润表处理.py
python roe_step3_TTM和平均权益.py
python roe_step4_点时间对齐_计算ROE.py
python roe_step5_MAD去极值_Z标准化.py
python verify_roe.py
\end{verbatim}

\subsection{Size}

\begin{verbatim}
python size_step1_复用EP行情.py
python size_step2_计算对数市值.py
python size_step3_MAD去极值_Z标准化.py
python verify_size.py
\end{verbatim}

\section{主要限制}

\begin{enumerate}
  \item 历史 ST/PT 标记缺失，当前结果不能声称已经完成历史 ST/PT 过滤；
  \item 官方上市日期缺失，次新股使用行情表首次出现日期作为代理；
  \item 2020 年初部分财报历史覆盖不足，早期因子覆盖率可能偏低；
  \item 只有因子计算，没有实际交易执行、涨跌停、成交量、滑点和交易成本模拟；
  \item \code{VERSION\_TIME} 使用 \code{max(ACT\_PUBTIME, UPDATE\_TIME)} 处理源表修订；若供应商能够提供更准确的修订公告时间，应优先使用官方公告时间；
  \item ROE 对非正平均归母权益设为缺失；PB 对负归母权益保留负 PB。两者处理不同，是因为负净资产对 PB 是有信息的，而 ROE 的非正权益分母会破坏收益率解释。
\end{enumerate}

\section{文件占用}

四个交付目录中的 Parquet 文件压缩后合计约：

\begin{table}[htbp]
\centering
\caption{四个交付目录的磁盘占用}
\begin{tabular}{lr}
\toprule
因子 & 占用\\
\midrule
EP & 257.11 MiB\\
PB & 258.35 MiB\\
ROE & 302.07 MiB\\
Size & 373.37 MiB\\
合计 & 约 1,190.9 MiB，即约 1.16 GiB\\
\bottomrule
\end{tabular}
\end{table}

上述数值是 Parquet 压缩后的磁盘占用，不等于 pandas 运行时的内存占用。

\end{document}

